% STATUS: written 2026-11-05 (Lopatkin 2019 read in full)
% Sources read in full: Lopatkin 2019 (pp. 59--89, sections 1 to 3).
% Origins: the Leavitt path algebra and its relations, Lopatkin 2019 (Definition 1.1, p. 61);
% the matrix and Laurent examples (Examples 1.2 and 1.3, p. 61); the basis
% (Theorem 1.6, p. 63); the form of a derivation (Theorem 2.1, p. 63); its
% generating set (Corollary 2.1, p. 64); the generators of the outer derivations
% (Theorem 3.2, p. 78); the dimension (Corollary 3.8, p. 79).
\section{Leavitt path algebras}\label{sec:path-incidence:leavitt}

The trivial extension of Section~\ref{sec:path-incidence:trivial-extension} adds one
crossed term to an algebra. The Leavitt path algebra of this section adds a reverse
edge for every arrow, and the presentation lists the reverse edges as new generators
\cite[p.~61]{Lopatkin2019}. Lopatkin describes every derivation of it on a row-finite
graph \cite[p.~59]{Lopatkin2019}. This section gives his description, the generating
set, and the dimension of the outer derivations.

The derivations below act on a Leavitt path algebra, the graded algebra the path and incidence cases of this part lead to.

A row-finite graph has finitely many edges with each source
\cite[p.~60]{Lopatkin2019}. The algebra has a vertex for each point of the graph, an
edge for each arrow, and a reverse edge for each arrow \cite[p.~61]{Lopatkin2019}.
Four sets of relations present it, the last of which splits a vertex among the edges
that leave it \cite[p.~61]{Lopatkin2019}.

\begin{definition}[{\cite[Definition~1.1, p.~61]{Lopatkin2019}}]\label{def:leavitt:algebra}
Let $K$ be a commutative ring and $\Gamma=(V,E)$ a row-finite graph. The
\emph{Leavitt path algebra} $L_{K}(\Gamma)$ of $\Gamma$ with coefficients in $K$ is
the $K$-algebra presented by the generators $\{v,\ v\in V\}$ and
$\{e,e^{*}\colon e\in E\}$ and the relations:
\begin{enumerate}[label=(\arabic*)]
\item $vu=\delta_{v,u}v$ for all $v,u\in V$;
\item $s(e)e=er(e)=e$ and $r(e)e^{*}=e^{*}s(e)=e^{*}$ for all $e\in E$;
\item $e^{*}f=\delta_{e,f}r(e)$ for all $e,f\in E$;
\item $v=\sum_{s(e)=v}ee^{*}$ for an arbitrary vertex $v\in V\setminus\{\text{sinks}\}$.
\end{enumerate}
\end{definition}

Two familiar algebras appear as examples \cite[p.~61]{Lopatkin2019}. A line of $n$
vertices gives the full matrix algebra \cite[p.~61]{Lopatkin2019}. A single loop
gives the Laurent polynomial algebra \cite[p.~61]{Lopatkin2019}.

\begin{example}[{\cite[Example~1.2, p.~61]{Lopatkin2019}}]\label{ex:leavitt:matrix}
Let $A_{n}$ denote the graph
\begin{equation*}
\bullet_{v_{1}}\xrightarrow{e_{1}}\bullet_{v_{2}}\xrightarrow{e_{2}}\cdots
\xrightarrow{e_{n-1}}\bullet_{v_{n}} .
\end{equation*}
Then $L_{K}(A_{n})\cong M_{n}(K)$, the full $n\times n$ matrix algebra. The
isomorphism sends $v_{i}$ to $E_{i,i}$, $e_{i}$ to $E_{i,i+1}$ and $e^{*}_{i}$ to
$E_{i+1,i}$ for $1\le i\le n-1$.
\end{example}

The second example has one vertex and one loop, and there the reverse edge acts as
the inverse variable \cite[p.~61]{Lopatkin2019}.

\begin{example}[{\cite[Example~1.3, p.~61]{Lopatkin2019}}]\label{ex:leavitt:laurent}
Let $R_{1}$ denote the graph with one vertex $v$ and one loop $e$. Then
$L_{K}(R_{1})\cong K[t,t^{-1}]$, the Laurent polynomial algebra. The isomorphism is
$v\mapsto 1$, $e\mapsto t$ and $e^{*}\mapsto t^{-1}$.
\end{example}

The algebra has a basis in normal form \cite[p.~63]{Lopatkin2019}. A word in the
basis is a path, a reverse path, or a pair of paths of equal length that end at the
same vertex \cite[p.~63]{Lopatkin2019}. The last pair must not both end in the same
special edge (Theorem~\ref{thm:leavitt:basis}) \cite[p.~63]{Lopatkin2019}.

The basis is what makes a derivation of $L_K(\Gamma)$ computable: a derivation is
determined by its values on a set of algebra generators, and the normal form
below reduces the product rule to a finite check on the words listed. Every
derivation statement in this section is read against that basis.

\begin{theorem}[{\cite[Theorem~1.6, p.~63]{Lopatkin2019}}]\label{thm:leavitt:basis}
Let $\Gamma$ be a row-finite graph and $K$ a commutative ring. Then the following
elements form a $K$-basis $\mathcal B$ of the Leavitt path algebra $L(\Gamma)$:
(1) the set of all vertices $V$; (2) the set of all paths $P$; (3) the set
$P^{*}=\{p^{*}\colon p\in P\}$; (4) a set $M$ of words of the form $wh^{*}$, where
$w=w_{0}\cdots w_{z}\in P$ and $h^{*}=[h_{0}\cdots h_{z}]^{*}=h^{*}_{z}\cdots h^{*}_{0}\in P^{*}$
are paths that end at the same vertex $r(w)=r(h)$, with the condition that the edges
$w_{z}$ and $h_{z}$ are either distinct or equal but not special.
\end{theorem}

Every derivation is fixed by its values on the generators
\cite[pp.~63--64]{Lopatkin2019}. Those values are sums over the basis, and the
coefficients satisfy three linear relations \cite[pp.~63--64]{Lopatkin2019}.

\begin{theorem}[{\cite[Theorem~2.1, pp.~63--64]{Lopatkin2019}}]\label{thm:leavitt:form}
Let $\Gamma$ be a row-finite graph and $K$ a commutative ring. Every derivation $D$
of the Leavitt path algebra $L_{K}(\Gamma)$ can be described as follows:
\begin{equation*}
D(v)=\sum_{wh^{*}\in\mathcal B}F_{wh^{*}}(v)\cdot\mathrm{ad}_{wh^{*}}(v),
\end{equation*}
\begin{equation*}
D(e)=\sum_{wh^{*}\in\mathcal B}F_{wh^{*}}(s(e))\mathrm{ad}_{wh^{*}}(e)
+\sum_{\substack{wh^{*}\in\mathcal B\\ s(w)=s(e),\ s(h)=r(e)}}F_{wh^{*}}(e)wh^{*},
\end{equation*}
\begin{equation*}
D(e^{*})=\sum_{wh^{*}\in\mathcal B}F_{wh^{*}}(s(e))\mathrm{ad}_{wh^{*}}(e^{*})
+\sum_{\substack{wh^{*}\in\mathcal B\\ s(w)=r(e),\ s(h)=s(e)}}F_{wh^{*}}(e^{*})wh^{*},
\end{equation*}
for every $v\in V$ and $e\in E$. Further, almost all coefficients are zero, and they
satisfy the relations
\begin{equation*}
F_{p}(e^{*})+F_{pf\!f^{*}}(e^{*})+F_{epf}(f)=0,\qquad epf\neq 0,
\end{equation*}
\begin{equation*}
F_{[f\!pe]^{*}}(e^{*})+F_{p^{*}}(f)+F_{e[pe]^{*}}(f)=0,\qquad f\!pe\neq 0,
\end{equation*}
\begin{equation*}
F_{w[fh]^{*}}(e^{*})+F_{ewh^{*}}(f)=0,\qquad w[fh]^{*}=pf\!f^{*},\ ewh=e[pe]^{*},
\end{equation*}
where $p,w,h\in P\cup V$ and $e,f\in E$ with $s(e)=s(f)$.
\end{theorem}

The derivations of Theorem~\ref{thm:leavitt:form} are generated by four families
\cite[p.~64]{Lopatkin2019}. Each family is attached to a path or a pair of paths (Corollary~\ref{cor:leavitt:generators})
\cite[p.~64]{Lopatkin2019}.\begin{corollary}[{\cite[Corollary~2.1, p.~64]{Lopatkin2019}}]\label{cor:leavitt:generators}
The $K$-module $\mathrm{Der}(L_{K}(\Gamma))$ is generated by the following set of
derivations:
\begin{enumerate}[label=(\arabic*)]
\item $\bigcup_{pq^{*}\in\mathcal B}\{\mathrm{ad}_{pq^{*}}\}$;
\item $\bigcup_{wh^{*}\in M,\ e\in E,\ s(w)=s(h)}\{D_{wh^{*}},D_{ee^{*}}\}$, where
$D_{wh^{*}}(h_{0})=w[h/h_{0}]^{*}$, $D_{wh^{*}}(w_{0}^{*})=-[w/w_{0}]h^{*}$, and
$D_{wh^{*}}=0$ otherwise;
\item $\bigcup_{c\in P,\ s(c)=r(c)}\{D_{c}\}$, where $D_{c}(e)=ce$ for $e\in E$,
$D_{c}(c_{0}^{*})=-[c/c_{0}]^{*}$, and $D_{c}=0$ otherwise;
\item $\bigcup_{c\in P,\ s(c)=r(c)}\{D_{c^{*}}\}$, where $D_{c^{*}}(e^{*})=-[ce]^{*}$
for $e\in E$, $D_{c^{*}}(c_{0}^{*})=[c/c_{0}]^{*}$, and $D_{c^{*}}=0$ otherwise.
\end{enumerate}
\end{corollary}

Modulo the inner derivations, the same families generate the outer derivations, with
explicit relations among them \cite[p.~78]{Lopatkin2019}. The relations record when a
generator dies or splits (Theorem~\ref{thm:leavitt:outer}) \cite[p.~78]{Lopatkin2019}.

\begin{theorem}[{\cite[Theorem~3.2, p.~78]{Lopatkin2019}}]\label{thm:leavitt:outer}
Let $\Gamma=(V,E)$ be a row-finite graph and $L_{K}(\Gamma)$ the Leavitt path
algebra over a commutative ring $K$. The $K$-module $HH^{1}(L_{K}(\Gamma))$ of outer
derivations of $L_{K}(\Gamma)$ can be presented by the following sets of generators:
\begin{enumerate}[label=(\arabic*)]
\item $\bigcup_{wh^{*}\in M,\ e\in E,\ s(w)=s(h)}\{D_{wh^{*}},D_{ee^{*}}\}$, where
$D_{ee^{*}}(e)=e$, $D_{ee^{*}}(e^{*})=-e^{*}$, and $D_{ee^{*}}=0$ otherwise, and
$D_{wh^{*}}(h_{0})=w[h/h_{0}]^{*}$, $D_{wh^{*}}(w_{0}^{*})=-[w/w_{0}]h^{*}$, and
$D_{wh^{*}}=0$ otherwise;
\item $\bigcup_{c\in P,\ s(c)=r(c)}\{D_{c}\}$, where $D_{c}(e)=ce$ for every $e\in E$,
$D_{c}(c_{0}^{*})=-[c/c_{0}]^{*}$, and $D_{c}=0$ otherwise;
\item $\bigcup_{c\in P,\ s(c)=r(c)}\{D_{c^{*}}\}$, where $D_{c^{*}}(e^{*})=-[ce]^{*}$
for every $e\in E$, $D_{c^{*}}(c_{0}^{*})=[c/c_{0}]^{*}$, and $D_{c^{*}}=0$ otherwise;
\end{enumerate}
and the following sets of relations among those generators:
\begin{enumerate}[label=(\arabic*)]
\item $D_{wh^{*}}=0$ whenever $|\{p\in P\colon r(p)=s(w)\}|<\infty$;
\item $D_{wh^{*}}=\sum_{e\in E}D_{ew[eh]^{*}}$ whenever
$|\{e\in E\colon r(e)=s(w)\}|<\infty$;
\item if the edge $e$ is special, $s(e)=r(e)$,
$\mathrm{supp}_{e}(s(e))\cap\mathrm{supp}_{e}(r(e))=\varnothing$, and both supports
are finite, then
\begin{equation*}
D_{ee^{*}}=\alpha\!\!\sum_{\substack{f,g\in E\setminus\vartheta(V)\\ r(g)\in\mathrm{supp}_{e}(s(e)),\ s(f)\in\mathrm{supp}_{e}(s(e))}}\!\!(D_{gg^{*}}-D_{f\!f^{*}})
+\beta\!\!\sum_{\substack{f,g\in E\setminus\vartheta(V)\\ r(g)\in\mathrm{supp}_{e}(r(e)),\ s(f)\in\mathrm{supp}_{e}(r(e))}}\!\!(D_{gg^{*}}-D_{f\!f^{*}}),
\end{equation*}
with $\alpha,\beta\in K$ and $\alpha-\beta=1$;
\item $D_{c}=D_{\zeta^{k}c}-\sum_{e\in E}D_{\zeta^{k}cee^{*}}$ whenever
$\delta_{c_{0},\vartheta}\cdots\delta_{c_{k-1},\vartheta}=0$,
$\delta_{c_{k},\vartheta}=1$ and $|\{p\in P\colon r(p)=s(c_{0}),\ p_{z}=c_{k}\}|<\infty$.
\end{enumerate}
\end{theorem}

Over a field the dimension of the space of outer derivations has two possible values
\cite[p.~79]{Lopatkin2019}. It is infinite exactly when the graph has an infinite
path \cite[p.~79]{Lopatkin2019}.

\begin{corollary}[{\cite[Corollary~3.8, p.~79]{Lopatkin2019}}]\label{cor:leavitt:dimension}
Let $K$ be a field. For the vector space $HH^{1}(L_{K}(\Gamma))$ of outer
derivations of the Leavitt path algebra $L_{K}(\Gamma)$ we have
\begin{equation*}
\dim HH^{1}(L_{K}(\Gamma))=
\begin{cases}
\infty, & \text{if there exists a path in }\Gamma\text{ with infinite length},\\
0, & \text{otherwise}.
\end{cases}
\end{equation*}
\end{corollary}

By Corollary~\ref{cor:leavitt:dimension}, a Leavitt path algebra carries as many
outer derivations as its graph has infinite paths. The record of the chapter is in
Sections~\ref{sec:path-incidence:path} to this one, and
Chapter~\ref{ch:polynomial-weyl} turns to algebras built from a polynomial ring.

Further work in the Leavitt and graded-ring line, held in the collection:
\noindent\emph{Prime group graded rings with applications to partial crossed products and Leavitt path algebras}~\cite{LibPrimeGroupGradedRings2021}.
